\chapter{アテンションとフィルター効果}
\section{カクテルパーティー効果と網様体賦活系（RAS）}
\epigraph{``In the music-room not a single light turned on... yet from out of the noise and the laughter came the name, whispered and repeated, leaping out above the roar of the party.''\\
（音楽室にはひとつの明かりも灯っていなかったが……喧騒と哄笑の中からその名前が立ち現れ、囁かれ、反復され、パーティーの轟音を突き破って飛び込んできた。）}{--- F. Scott Fitzgerald, \textit{The Great Gatsby}}

騒がしい立食パーティーや雑踏の喧騒の中にいる場面を想像してみてほしい。
周囲にはグラスがぶつかる音、BGM、無数の人々の話し声が混ざり合い、物理的には強大な音響ノイズの塊（$a$）が鼓膜を激しく叩き続けている。
通常、それらの雑音は意識されることなく切り捨てられ、目の前の相手との会話にだけ焦点が合っている。

ところが、その喧騒のどこか遠くで、誰かが不意に「あなたの名前」を口にした瞬間、どうなるだろうか。
どんなに周囲のノイズ音圧が高く、その声自体が小さかったとしても、意識は瞬時にそれを捉え、身体は反射的にその声の主の方へと振り向くはずだ。
心理学において「カクテルパーティー効果」として知られるこの現象は、私たちの脳が外界の音波を均一に記録するマイクではなく、特定の信号を選択的に増幅・遮断する物理フィルターを内蔵していることを証明している。

この生物学的ハードウェアの正体が、脳幹から大脳皮質へと投射される神経ネットワーク、すなわち\textbf{網様体賦活系（RAS: Reticular Activating System / 上行性網様体賦活系）}である。

RASは脳幹（中脳・橋・延髄）の網様体を起点とし、間脳（視床・視床下部）を経由して大脳皮質全体へと広く神経線維を伸ばしている。
その主要な機能は、意識の覚醒度をコントロールすることに加え、感覚器から絶え間なく流入する莫大な入力情報に対して「重要度に応じたゲートキーパー（選別機）」として機能することにある。

この機構を神経系疑似ベイズ推論機のモデルに組み込むと、その物理的意味が明瞭になる。
感覚器を叩く入力ノイズ $e_{\mathrm{noise}}$ の情報量は毎秒数百メガビットにも達するが、意識（大脳皮質）が直列処理できる帯域はわずか数十ビット程度にすぎない。
もしすべてのノイズを均等に計算処理しようとすれば、脳の計算資源は瞬時にパンクする。

そこでRASは、過去の記憶や情動と結びついた「自己の名前」や「危険シグナル」といった特定のパターンに対して、極めて高い事前予期 $\hat{\omega}$ のバイアスをあらかじめセットしておく。
雑音の海からその特徴量（名前の音素配列など）が微小でも検出された瞬間、RASは皮質へ覚醒信号を送り、その信号の尤度を一気に引き上げて閾値を突破させる。
\begin{equation}
    a_{\mathrm{noise}} \nbuc{\hat{\omega}_{\mathrm{name}}}{\hat{\omega'}_{\mathrm{name}}} \omega_{\mathrm{name}} \equiv A_e
\end{equation}

名前を呼ばれて振り返るという日常的な反応は、意識が能動的に聞き耳を立てた結果ではない。
RASという神経系の物理ゲートが、無意識下でノイズの海をスクリーニングし、事前設定されたパターンに合致した瞬間に確信 $\omega$ を強制的に確定させ、身体の反射運動をトリガーした結果なのである。


